\chapter{Trabalhos futuros}

O EIGAM chega ao fim deste trabalho em beta fechado, com a arquitetura pronta para crescer e com um conjunto de lacunas conhecidas. Os trabalhos futuros se dividem em quatro frentes.

\textbf{Validação com usuários reais.} O passo mais importante é abrir o produto para um grupo maior e medir o que este trabalho não pôde medir: satisfação, retenção, frequência de uso, conversão do período de teste em assinatura e mudança no comportamento financeiro. O mecanismo de avaliação por nota e comentário já existe; falta volume. Com ele, será possível aplicar um questionário de usabilidade, revisar amostras de conversas por critério de risco e calibrar os limiares dos avisos com dados de muitos titulares, e não apenas de alguns. Se a conversão do teste sem cartão se mostrar baixa, ou o custo de conexão bancária de quem não converte se mostrar alto, o teste com cartão registrado é a alternativa documentada. Antes da abertura, a política de privacidade precisa ser atualizada para refletir a leitura de investimentos e empréstimos, já implementada, e a integração com assistentes externos depende das revisões jurídica e de segurança dedicadas, incluindo o tratamento da transferência internacional de dados quando os valores são liberados.

\textbf{Previsão e avisos.} As projeções estão implementadas e medidas; o que falta é o veredito com dados reais. Os próximos passos são acumular fechamentos mensais da coorte do beta no registro de previsões, confirmar ou descartar o candidato em sombra que decompõe o mês em compromissos conhecidos e gasto variável, reavaliar a faixa do saldo livre com doze meses de histórico, cuja cobertura ainda está abaixo da régua, e só então ligar as chaves de liberação para o público geral. A jornada completa de metas e alertas em canal real, último portão da liberação, inclui duas semanas de acompanhamento em sombra dos avisos de orçamento. Os alertas de vencimento e de saldo insuficiente, previstos na proposta inicial, podem ser construídos sobre a detecção de recorrências, que já identifica a próxima cobrança, sobre a posição de caixa do saldo livre e sobre a mesma política de avisos usada pelos demais sinais.

\textbf{Qualidade e desempenho.} Os \textit{golden sets} devem passar a bloquear a integração contínua, como a arquitetura previa, e as baterias com o modelo real podem ser agendadas em vez de executadas manualmente. A latência das respostas complexas pode ser reduzida com um modelo de redação mais rápido, com \textit{cache} semântico para perguntas recorrentes e com respostas tipadas montadas pelo servidor em mais tipos de consulta. A troca do cliente Kafka do \textit{Chat Service} por um que não espere o lote inteiro tornaria desnecessário o arranjo de um consumidor por partição.

\textbf{Operação e produto.} A infraestrutura deve ser separada em servidores dedicados para aplicação, dados e \textit{broker}, com estado remoto do Terraform, rotação automatizada de chaves e réplica de leitura do banco a partir de alguns milhares de usuários. A migração para servidores ARM reduz o custo em cerca de 20\%. No produto, a publicação dos \textit{WhatsApp Flows}, que aguarda a liberação do provedor do canal, levará ao \textit{WhatsApp} os mesmos formulários que já funcionam no \textit{Telegram}, e a arquitetura multicanal permite avaliar novos canais pelo custo e pelo alcance de cada um. A integração com assistentes externos pode evoluir com a publicação do conector nos diretórios dos provedores de assistentes, que exigem política de privacidade e documentação públicas, e com o suporte a identificação de clientes por documentos de metadados, caso o volume de conexões cresça. Por fim, a orientação sobre o próprio orçamento pode evoluir para planos de ação acompanhados ao longo do mês, sempre dentro do limite que separa educação financeira de recomendação de investimento.
